\section{探索性分析与数据可视化}

本章的目的在于通过可视化理解样本分布、变量相关性和不同群体之间的描述性差异。EDA 不直接证明因果关系，但可以帮助确定模型任务的难点，例如类别不均衡、变量相关性和风险组之间的行为差异。

\subsection{高风险标签分布}

图 \ref{fig:eda-high-risk-distribution} 展示了 \texttt{high\_risk\_flag} 的类别分布。结果文件 \texttt{results/eda\_high\_risk\_flag\_distribution.csv} 显示，低风险样本为 2795 条，占 79.86\%；高风险样本为 705 条，占 20.14\%。这说明分类任务存在一定类别不均衡，如果只看 Accuracy，模型可能偏向多数类。因此后续分类实验需要同时报告 Recall、F1、PR-AUC 和混淆矩阵。

\maybefigure{eda_high_risk_flag_distribution.png}{\texttt{high\_risk\_flag} 类别分布}{fig:eda-high-risk-distribution}

\subsection{数值变量相关性}

图 \ref{fig:eda-correlation} 展示主要数值变量之间的相关性结构。该图用于观察数字行为、睡眠活动和结果变量之间是否存在较强线性关联，也用于提醒后续模型解释时注意变量之间可能存在共线性。相关性图只能说明变量之间的统计关联强弱，不能说明设备使用、睡眠或通知数量对结果变量存在因果作用。

\maybefigure{eda_numeric_correlation_heatmap.png}{数值变量相关性热力图}{fig:eda-correlation}

\subsection{风险组行为差异}

图 \ref{fig:eda-boxplots-risk} 按 \texttt{high\_risk\_flag} 对核心行为变量进行箱线图比较。根据 \texttt{results/eda\_group\_comparison\_summary.csv}，高风险组的日均设备使用时长均值为 9.47 小时，非高风险组为 6.77 小时；高风险组手机解锁次数均值为 186.92 次，非高风险组为 137.06 次；高风险组睡眠时长均值为 6.50 小时，非高风险组为 7.45 小时。这些差异为分类模型提供了可解释的行为线索，但仍应理解为描述性差异。

\maybefigure{eda_boxplots_by_risk.png}{按高风险标签分组的行为变量箱线图}{fig:eda-boxplots-risk}

\subsection{类别变量风险率}

图 \ref{fig:eda-category-risk-rate} 汇总了性别、地区、收入水平、教育水平、日常角色和设备类型等类别变量下的高风险比例。该图的作用不是给某个群体贴标签，而是检查类别变量是否可能对模型产生分组差异。由于数据为合成数据，类别风险率仅用于描述样本结构和辅助建模解释，不用于真实人群判断。

\maybefigure{eda_category_risk_rate.png}{类别变量分组风险率}{fig:eda-category-risk-rate}
